\documentclass[11pt]{article}

\usepackage[margin=1in]{geometry}
\usepackage{amsmath, amssymb}
\usepackage{mathrsfs}
\usepackage{bm}
\usepackage{xcolor}
\usepackage{fancyhdr}
\usepackage{booktabs}
\usepackage{multirow}
\usepackage{array}
\usepackage{url}
\usepackage{soul}
\usepackage{pdfcomment}
\definecolor{hlyellow}{RGB}{255,255,170}
\definecolor{hlcyan}{RGB}{170,255,255}
\definecolor{hlgreen}{RGB}{170,255,170}
\definecolor{hlpink}{RGB}{255,170,170}
\definecolor{hlorange}{RGB}{255,221,170}

\pagestyle{plain}


\begin{document}

% =========== Page 1: ergodic_page1.tex ===========
\setcounter{equation}{0}

% ---------- IEEE acceptance banner ----------
\begin{center}
\footnotesize
\pdfcomment[color=green,subject={},open=false,opacity=0.7]{User input: that paper was PUBLISHED after Fall 2025 in ergodic exploration paper\\Claude comments: EVIDENTIAL RELATIONSHIP: supports. DOI embeds 2026; masthead states ACCEPTED FEBRUARY, 2026. Fall 2025 ended by mid-December 2025; Feb 2026 is $\sim$2 months after.}\sethlcolor{hlgreen}\hl{This article has been accepted for publication in IEEE Robotics and
Automation Letters. This is the author's version which has not been
fully edited and content may change prior to final publication.
Citation information: DOI 10.1109/LRA.2026.3673907}

\vspace{0.3em}
\textbf{\sethlcolor{hlgreen}\hl{IEEE ROBOTICS AND AUTOMATION LETTERS. PREPRINT VERSION.
ACCEPTED FEBRUARY, 2026}}
\end{center}

\vspace{1em}

% ---------- Title block ----------
\begin{center}
{\Large\bfseries Autonomous Robotic Tissue Palpation and Abnormalities
Characterisation via Ergodic Exploration}

\vspace{0.8em}

{\normalsize Luca Beber$^{1}$, Edoardo Lamon$^{1}$,
Matteo Saveriano$^{2}$, Daniele Fontanelli$^{2}$, and
Luigi Palopoli$^{1}$}
\end{center}

\vspace{0.8em}

% ---------- Abstract ----------
\noindent\textbf{Abstract}---We propose a novel autonomous robotic
palpation framework for real-time elastic mapping during tissue
exploration using a viscoelastic tissue model. The method combines
force-based parameter estimation using a commercial force/torque
sensor with an ergodic control strategy driven by a tailored Expected
Information Density, which explicitly biases exploration toward
diagnostically relevant regions by jointly considering model
uncertainty, stiffness magnitude, and spatial gradients. An Extended
Kalman Filter is employed to estimate viscoelastic model parameters
online, while Gaussian Process Regression provides spatial modelling
of the estimated elasticity, and a Heat Equation Driven Area Coverage
controller enables adaptive, continuous trajectory planning.
Simulations on synthetic stiffness maps demonstrate that the proposed
approach achieves better reconstruction accuracy, enhanced
segmentation capability, and improved robustness in detecting stiff
inclusions compared to Bayesian Optimisation-based techniques.
Experimental validation on a silicone phantom with embedded inclusions
emulating pathological tissue regions further corroborates the
potential of the method for autonomous tissue characterisation in
diagnostic and screening applications.

\vspace{0.6em}

\noindent\textbf{Index Terms}---Medical Robots and Systems;
Sensor-based Control; Force and Tactile Sensing.

\vspace{0.8em}

% ---------- Fig. 1 placeholder ----------
\begin{center}
\fbox{\begin{minipage}[c][6cm][c]{0.90\textwidth}
\centering
\textbf{Fig. 1.}\ Experimental setup for the proposed ergodic search.\\[4pt]
\textcolor{gray}{\footnotesize
\textbf{(bottom-left)} silicone sample (Ecoflex-0030 matrix with a
DragonSkin-30NV spherical inclusion);\\
\textbf{(top-left)} ground-truth elasticity map (Elastic Modulus in kPa)
obtained using a grid-based approach;\\
\textbf{(right)} the UR3e robotic manipulator palpating the silicone
sample. Labelled components: F/T Sensor (Bota SensOne, between flange
and indenter) and Indenter (3D-printed spherical tip, 5 mm radius).}
\end{minipage}}
\end{center}

\vspace{0.4em}

\textbf{Fig. 1.}\ Experimental setup for the proposed ergodic search:
(bottom-left) silicone sample; (top-left) ground-truth elasticity map
obtained using a grid-based approach; (right) the robotic manipulator
palpating the silicone sample.

\vspace{1em}

% ---------- Section I ----------
\section*{I.\ Introduction}

\textbf{P}ALPATION is a fundamental clinical procedure that allows
physicians, by touch, to assess tissue stiffness, texture, and the
presence of abnormalities such as tumours or swollen lymph nodes [1].
Its speed and non-invasiveness make it crucial for early diagnosis,
particularly of conditions such as breast cancer. However, palpation
is inherently subjective and sensitive to inter-practitioner
variability, which may compromise diagnostic accuracy and
repeatability for small or deep-seated lesions. For this reason,
physical examinations are often complemented or replaced by imaging
techniques such as mammography, the gold standard for early detection
[2], and ultrasound [3].

Accurate assessment of tissue mechanics is equally critical in
robot-assisted minimally invasive surgery (RMIS), where robots must
adapt to variable stiffness for safe navigation and manipulation,
especially for autonomous systems operating with minimal human
supervision [4]. To support such needs, several technologies have
emerged. Imaging-based elastography methods, e.g., ultrasound [5],
magnetic resonance [6], or laser-based methods [7], estimate tissue
stiffness but typically provide relative measures, are costly or slow,
and cannot be easily integrated into RMIS. Tactile sensing, in
contrast, offers real-time, quantitative stiffness measurements by
capturing force and deformation. Previous works have shown its
potential to reduce applied forces, shorten procedures, and improve
surgical outcomes [8], [9]. Early stiffness mapping used uniform
probing grids [10], [11], while later methods adopted Bayesian
Optimisation (BO) to reduce the number of probing actions [12]--[14],
often coupled with custom tactile sensors [8], [9], [15], guided by
anatomical priors [16] or with ad-hoc refined segmentation strategies
[15], [17].

In this work, we propose a method for autonomous robotic palpation
that enables continuous estimation of elastic tissue stiffness using
a viscoelastic contact model, relying solely on a commercial
force/torque (F/T) sensor. The use of a viscoelastic model allows
more accurate representation of tissue dynamics during palpation
compared to purely elastic formulations, while abnormality
characterisation in this work focuses on the estimated elastic
stiffness. \pdfcomment[color=orange,subject={},open=false,opacity=0.7]{User input: that the Ergodic exploration paper has instrumentation that is designed for EASY ADOPTION meaning that people bad at experiments can use it and so can people just use it at home - that the barrier for adoption is LOW\\Claude comments: EVIDENTIAL RELATIONSHIP: partially\_supports. Low-barrier thesis strongly supported (off-the-shelf, portability, cost-effectiveness, broader applicability, standard-platform compatibility). But 'home use' and 'non-experts' are NOT claimed by the paper; target audience is robotics/RMIS research community.}\sethlcolor{hlorange}\hl{This sensing strategy provides a practical alternative to
methods based on embedded or specialised sensors, which typically
require customised hardware integration and lack compatibility with
standard robotic platforms. In contrast, the use of an off-the-}

\vspace{1em}

% ---------- Footer: submission + affiliations ----------
\noindent\rule{0.4\textwidth}{0.4pt}\\[4pt]
{\footnotesize
Manuscript received: October, 15, 2025; Revised December, 16, 2025;
Accepted February, 23, 2026.\\[2pt]
This paper was recommended for publication by Editor Burgner-Kahrs
Jessica upon evaluation of the Associate Editor and Reviewers'
comments. This research has been funded by the the PNRR project
FAIR - Future AI Research (PE00000013), by the European Union
projects INVERSE (GA no.~101136067) and MAGICIAN (GA no.~101120731)
and by the MUR Departments of Excellence 2023-27 program
(L.232/2016).\\[2pt]
$^{1}$ Department of Information Engineering and Computer Science,
University of Trento, Trento, Italy. \texttt{name.surname@unitn.it}\\[2pt]
$^{2}$ Department of Industrial Engineering, University of Trento,
Trento, Italy. \texttt{name.surname@unitn.org}\\[2pt]
Digital Object Identifier (DOI): see top of this page.\\[6pt]
This work is licensed under a Creative Commons Attribution 4.0
License. For more information, see
\texttt{https://creativecommons.org/licenses/by/4.0/}
}

\newpage

% =========== Page 2: ergodic_page2.tex ===========
\setcounter{equation}{0}

% ---------- Continuation of Section I Introduction ----------
\sethlcolor{hlorange}\hl{shelf F/T sensor ensures portability, cost-effectiveness, and broader
applicability of the proposed framework.} We couple a force-based
contact model implemented via the Dimensionality Reduction Method
(DRM) [18] with an ergodic exploration strategy to efficiently
construct stiffness maps, where the robot explores the environment
in proportion to the Expected Information Density (EID). To plan
trajectories, we employ the Heat Equation Driven Area Coverage
(HEDAC) algorithm [19], a real-time ergodic controller that balances
exploration and exploitation at 100 Hz. This replaces earlier ergodic
control methods, like Spectral Multiscale Coverage (SMC), which was
limited to 4 Hz as reported in [16]. The proposed framework is
evaluated in simulation and real-world experiments, including
validation on a silicone sample with a stiff inclusion (Fig.~1),
demonstrating robust performance and, in some cases, superior
detection and boundary delineation compared to BO.

To summarise, the main contributions of this paper are:
\begin{itemize}
\item The design of an EID tailored to autonomous stiffness mapping,
  which explicitly combines uncertainty, stiffness magnitude, and
  spatial gradients to guide ergodic exploration toward diagnostically
  relevant regions;
\item \pdfcomment[color=yellow,subject={},open=false,opacity=0.7]{User input: that the ergodic exploration paper is an example of adaptive sampling\\Claude comments: Cleanest statement of the paper's adaptive-sampling architecture. The phrase 'stiffness estimates continuously reshape the information objective driving exploration' is the feedback-loop signature of adaptive sampling: the sampling policy is continuously updated based on what has been learned, rather than fixed in advance. Alternatives: page 1 abstract (trajectory-focused), page 8 conclusion (mechanism implicit), page 3 Eq. 3 $\alpha(t)$ (mathematical but not prose). The page-2 bullet wins on being both self-contained and mechanism-explicit.}\sethlcolor{hlyellow}\hl{A closed-loop combination between online force-based
  viscoelastic parameter estimation and ergodic trajectory planning,
  where stiffness estimates continuously reshape the information
  objective driving exploration;}
\item A real-time implementation of ergodic palpation using HEDAC,
  enabling continuous exploration--exploitation trade-offs at 100 Hz
  with continuous motion;
\item A convergence-based stopping criterion derived from the
  ergodic metric.
\end{itemize}

% ---------- Section II ----------
\section*{II.\ Related Work}

Typically, search algorithms rely on three key components: (i) a
model of the underlying distribution, (ii) a strategy to select the
next sampling location, and (iii) a mechanism to store and update the
collected data. Although Gaussian Process Regression (GPR) [20] is
widely adopted to model spatially varying physical properties (ii),
such as tissue stiffness, the choice of sampling strategy (i) and
planning (iii) remains an active area of research.

\subsection*{A.\ Bayesian Optimisation for Palpation}

Several studies have employed BO to guide robotic palpation to reduce
the number of probing actions while maximising information gain. For
example, Garg et al.\ [12] investigated how different acquisition
functions affect the estimated stiffness map, showing how their
choice influences both peak detection and regional estimation. Yan
et al.\ [15] highlighted a key limitation of BO: while it excels at
locating stiffness peaks, it struggles to delineate their boundaries
accurately. To address this, they proposed a two-step approach
combining BO-based search and boundary refinement using a radial
strategy originating from the centroid of the estimated region.
However, this strategy requires a reliable estimation of the
centroid, which may not be guaranteed after exploration, and may be
less adaptive to irregular or non-convex boundaries. Expected
Improvement has been used as an acquisition function to improve
stiffness estimation [13], and later extended with a utility-guided
approach [16] that incorporates prior beliefs and penalises excessive
motion. Chalasani et al.\ proposed a method that employs GPR to
estimate tissue stiffness, which is then used as a prior to fit a
second GPR model for reconstructing the tissue surface geometry [21].
In a follow-up study, they extended this approach to support online
stiffness mapping during surgery [22]. A critical limitation of BO
is that exploration and exploitation are handled sequentially through
the choice of the acquisition function. As a result, BO may focus
excessively on promising regions, thereby overlooking smaller or
subtler areas that could yield informative measurements.

\subsection*{B.\ Ergodic Exploration}

Ergodic control has emerged as a compelling alternative to BO for
tasks that involve spatially distributed information, where
exploration must account for system dynamics. Unlike point-wise
optimisation strategies, which select discrete sampling locations,
ergodic control generates trajectories that allocate time across the
environment proportionally to the EID, making it well suited for
continuous spatial estimation tasks. Mathew and Mezi\'c [23]
introduced the Spectral Multiscale Coverage (SMC) framework, using
Fourier-based metrics to quantify how well a trajectory covers a
given spatial distribution. Miller et al.\ [24] demonstrated the use
of ergodic control for information-guided exploration, integrating
the Linear Quadratic Regulator (LQR) formulation of SMC [25].
However, the high computational cost of optimisation limited
real-time adaptation. Ayvali et al.\ [16] presented the first
application of ergodic control to palpation tasks using SMC,
concluding that BO achieved superior performance, particularly when
prior information was available. Conversely, our results indicate
the opposite trend, highlighting the potential of ergodic control
for robotic palpation. The main reason is that SMC often prioritises
global exploration and can suffer from over-sampling already explored
regions, while its reliance on spectral transforms poses challenges
for real-time deployment. To address these limitations, Ivic et al.\
[19] proposed HEDAC, which leverages radial basis functions and a
potential field derived from a stationary heat equation to compute
smooth local trajectories. This method improves scalability and
enables real-time implementation, avoiding the computational burden
of Fourier analysis.

% ---------- Section III opening ----------
\section*{III.\ Methodology}

In this work, we combine online force-based parameter estimation
using a viscoelastic tissue model with real-time ergodic control
implemented via HEDAC, enabling continuous robotic palpation without
tactile sensors and without sacrificing spatial coverage or
responsiveness. While ergodic control and online viscoelastic
estimation have been previously investigated, we show that effective
ergodic palpation critically depends on the definition of the target
distribution guiding exploration. We therefore propose an EID
tailored to stiffness mapping and lesion detection, which balances
exploration and exploitation by jointly considering uncertainty,
stiffness magnitude, and

\newpage

% =========== Page 3: ergodic_page3.tex ===========
\setcounter{equation}{0}

% ---------- Continuation of Section III opening ----------
spatial gradients. Although viscoelastic modelling supports robust
estimation, lesion search and characterisation in this work are
driven solely by elasticity. Our framework is structured as a
closed-loop search algorithm comprising four main components:
A.~a target distribution over tissue stiffness, which is updated at
every iteration; B.~an Expected Information Density (EID) map to
prioritise sampling locations; C.~a trajectory planner using ergodic
control, guided by the EID; D.~a data acquisition executed during
trajectory execution, which enables updates in the target
distribution. These steps are executed cyclically, updating the
target distribution (E.) until a predefined stopping criterion is
met (F.). Each component is designed to be modular and independent,
allowing flexible integration of alternative models or control
strategies.

\subsection*{A.\ Target Distribution}

\pdfcomment[color=green,subject={},open=false,opacity=0.7]{User input: that global maps are assembled by stitching together local estimates in the ergodic exploration paper\\Claude comments: EVIDENTIAL RELATIONSHIP: supports. GPR IS the stitching mechanism. 'Robot end-effector positions' = discrete local sampling sites; mean function $\mu(x)$ = stitched global stiffness map. The architectural pattern matches Chen's POINT$\to$MAP ladder precisely --- with GPR interpolation replacing direct ($x, y, z$)-tuple tabulation. Same local-estimates-to-global-map signature in both papers.}\sethlcolor{hlgreen}\hl{We use GPR to construct a continuous map of the stiffness
distribution across the palpated surface. The robot end-effector
positions serve as the input domain, and the estimated elasticity
values form the output. After each regression update, the mean
function $\mu(x)$ estimates the stiffness at unobserved locations,
while the variance function $\sigma^{2}(x)$ quantifies the model
uncertainty. As expected with GPR, uncertainty increases with
distance from the observed samples.} The squared exponential kernel
ensures smooth, differentiable estimates, which are critical for
downstream gradient computation.

\textbf{1) Gaussian Process Regression:} is a non-parametric Bayesian
approach for estimating continuous functions from data, offering
both predictions and uncertainty estimates. Unlike parametric
models, which assume a fixed functional form, GPR models a
distribution over functions that is updated as new observations are
incorporated [20]. A Gaussian Process (GP) is a collection of random
variables such that any finite subset follows a joint Gaussian
distribution. Formally, it is written as:
\[
  f(x) \sim \mathcal{GP}(m(x),\, k(x, x')),
\]
where $m(x)$ is the mean function (often assumed zero) and
$k(x, x')$ is the covariance function, or kernel, which encodes the
similarity between inputs. The kernel defines how closely related
the outputs are at different input points. One of the most widely
used kernels is the squared exponential (SE):
\[
  k_{\mathrm{SE}}(x, x') =
  \exp\!\left(-\frac{\|x - x'\|^{2}}{2 l^{2}}\right),
\]
where $l$ is the length scale parameter controlling smoothness
(in our experiments, $l = 2.5$ cm is equal to the radius of the
indenter), and $\|x - x'\|$ is the Euclidean distance between
inputs. This kernel implies smooth, infinitely differentiable
functions, making it ideal for modelling physical quantities such
as tissue stiffness.

\textbf{2) GPR Training:} Given a set of training observations $X$
with corresponding outputs $y$, GPR computes a posterior predictive
distribution at new input locations $X^{*}$. This yields a predicted
mean function $\mu^{*}$, representing the estimated stiffness at
unobserved locations, and an associated covariance $\Sigma^{*}$,
which quantifies model uncertainty. As expected, the uncertainty
increases with distance from the observed samples, reflecting
reduced confidence in sparsely explored regions.

\subsection*{B.\ Expected Information Density for Elasticity Maps}

To prioritise regions for further palpation, we define an EID
function $\xi_{\mathrm{EID}}(x)$ that balances two objectives:
exploration of uncertain areas, and exploitation of predicted
high-stiffness regions and refinement near region boundaries.
Formally, we define
\begin{equation}
\tilde{\xi}_{\mathrm{EID}}(x) =
(1 - \alpha)\,\frac{g(x)}{\int_{\mathcal{X}} g(x)}
+ \frac{\mu(x)}{\int_{\mathcal{X}} \mu(x)}
+ \alpha\,\frac{\sigma(x)}{\int_{\mathcal{X}} \sigma(x)}
\end{equation}
the un-normalised distribution, where
$g(x) = \|\nabla \mu(x)\|$ is the gradient norm of the predicted
elasticity, highlighting boundaries. The distribution (1) is then
normalised as:
\begin{equation}
\xi_{\mathrm{EID}}(x) =
\frac{\tilde{\xi}_{\mathrm{EID}}(x)}%
{\int_{\mathcal{X}} \tilde{\xi}_{\mathrm{EID}}(x)\, dx}.
\end{equation}
The trade-off between exploration and exploitation is controlled by
the scalar $\alpha \in [0, 1]$. Lower values bias the system towards
known high-stiffness or high-gradient areas, while higher values
encourage sampling in regions of high uncertainty. In our
implementation, $\alpha$ varies over time depending on the value of
the ergodic metric in the following way:
\begin{equation}
\alpha(t) =
\frac{\int_{\mathcal{X}} e(x, t)\, dx}%
{\int_{\mathcal{X}} e(x, 0)\, dx},
\end{equation}
where $e(x, t)$ is the ergodic metric that will be defined in the
following section in (4). According to (3), at the onset of the
exploration, the trajectory is not ergodic and $\alpha(0) = 1$. As a
consequence, the EID in (1) is dominated by the variance of the GPR,
resulting in a purely exploratory behaviour. Regions that have not
yet been sampled exhibit higher uncertainty and are therefore
prioritised for exploration. As informative measurements are
collected, the ergodic metric progressively decreases, leading to a
gradual rebalancing between exploration and exploitation. The
spatial coverage of the trajectory increasingly matches the target
EID distribution, and the exploration naturally shifts towards
exploitation of regions associated with higher estimated stiffness
and pronounced spatial gradients, such as the boundaries of stiff
inclusions. We want to highlight that the EID is computed on-the-fly
and does not exploit prior knowledge about the inspected area.

\subsection*{C.\ Trajectory Planning}

The normalised EID $\xi_{\mathrm{EID}}$ is used as the target
distribution for the ergodic planner, which computes a trajectory
that visits each region in proportion to its expected information
content. The HEDAC algorithm [19] is used to generate the trajectory
online.

\textbf{1) Ergodic Control:} aims to guide a robot trajectory to
match a target spatial distribution, balancing exploration and
exploitation by spending more time in regions of higher interest. It
models the coverage density using radial basis functions (RBFs).
Given a trajectory
$z: [0, t] \to \Omega \subset \mathbb{R}^{n}$,
the normalised coverage density $c(x, t)$ is:
\[
c(x, t) = \frac{\tilde{c}(x, t)}{\int_{\Omega} \tilde{c}(x, t)\,dx},
\qquad
\tilde{c}(x, t) = \frac{1}{t}\int_{0}^{t}
\phi_{\sigma}(x - z(\tau))\,d\tau,
\]

\newpage

% =========== Page 4: ergodic_page4.tex ===========
\setcounter{equation}{3}

% ---------- Continuation of Section III-C ergodic control defn ----------
where $\phi_{\sigma}$ is a Gaussian kernel. The coverage error with
respect to a target distribution $m(x)$ is defined as the spatial
diffusion of $m(x)$ through the RBF kernel compared with the coverage
density
\begin{equation}
e(x, t) = (\phi_{\sigma} * m)(x) - c(x, t),
\end{equation}
To avoid local minima and improve spatial coordination, HEDAC
smooths this error using a stationary heat equation:
\[
\varrho\, \Delta u(x, t) = \beta\, u(x, t) + \gamma\, a(x, t) - s(x, t),
\qquad
\dot{z} = v_{a} \cdot \frac{\nabla u}{\|\nabla u\|},
\]
where $s$ highlights under-sampled regions, and $a$ ensures
collision avoidance between the agents.

\textbf{2) Practical Implementation:} Different from [19], this work
implemented the non-stationary diffusion equation, as was done in
[26], for a more locally consistent exploration behaviour. This means
that $\dot{u} \ne 0$ is defined as:
\begin{equation}
\dot{u}(x, \tau) = \Delta u(x, \tau),
\end{equation}
allowing us to control the desired smoothness. In this formulation,
an initial condition is required: $u(x, 0) = s(x, 0)$. During the
experiments, a second-order agent is used to generate the trajectory.
The trajectory is updated with a frequency $\Delta T$ chosen to
ensure smooth motion and adequate control loop resolution (typically
100 Hz in our implementation, 5 times slower than the low-level
controller). The planning phase respects constraints on maximum
velocity to ensure physical plausibility during data collection
imposed by the second-order dynamics inside the ergodic controller.

% ---------- Section III-D ----------
\subsection*{D.\ Trajectory Execution and Data Collection}

While the trajectory is executed, tissue elasticity is estimated
continuously using an Extended Kalman Filter (EKF). Data are sampled
at a rate that depends on the robot motion and the indenter geometry.
Specifically, higher velocities or smaller indenters require denser
sampling to avoid aliasing effects in the stiffness map.

\textbf{1) Online Elasticity and Viscosity Estimation:} Accurately
estimating the mechanical response of soft tissues in real time is a
challenging task, as it requires modelling the contact interaction
between the probe and the tissue surface and inferring tissue
parameters from force measurements and robot motion. In this work,
we adopt a viscoelastic contact model to improve the fidelity of
force prediction during dynamic palpation, while subsequent
abnormality detection focuses on the estimated elastic stiffness.
To achieve this, we adopt DRM [27], which allows mapping the
three-dimensional contact mechanics into an equivalent
two-dimensional representation. \pdfcomment[color=green,subject={},open=false,opacity=0.7]{User input: that the quantities inferred from the measurements use a mathematical model that ONLY ACCOUNTS FOR THE LOCAL PHYSICS in the ergodic paper\\Claude comments: EVIDENTIAL RELATIONSHIP: supports. Hertzian contact mechanics + viscoelastic damping = purely local physics. Eq.~6 takes only local point-contact variables ($d, \dot{d}, E_f, \eta, \nu, R$) --- no bulk PDE, no global FEM solve, no neighbour coupling. The Hertzian half-space assumption is the definitive local-physics fingerprint. Contrast with MBT which solves a global elastic PDE.}\pdfcomment[color=green,subject={},open=false,opacity=0.7]{User input: that the inference model treats forces as NORMAL, i.e. no TANGENTIAL component\\Claude comments: EVIDENTIAL RELATIONSHIP: supports. $F_{TOT}$ is a scalar (not vector) --- no normal/tangential decomposition. All inputs ($d, \dot{d}, R, E_f, \eta, \nu$) are rotationally symmetric about the contact normal. No friction coefficient $\mu$ or shear modulus $G$ appears. The EKF state has no tangential kinematic variables. Hertzian contact mechanics is the canonical frictionless-normal-load theory.}\sethlcolor{hlgreen}\hl{Under this framework, the force
exerted by the tissue during indentation is a nonlinear function of
penetration depth and velocity:}
\begin{equation}
F_{\mathrm{TOT}}(d, \dot{d}) =
\frac{4}{3}\, \frac{E_{f}}{1 - \nu^{2}}\, \sqrt{R\, d}\, d
+ \frac{4}{1 - \nu}\, \eta\, \sqrt{R\, d}\, \dot{d},
\end{equation}
\sethlcolor{hlgreen}\hl{where $d$ is the penetration depth, $\dot{d}$ its derivative
(velocity), $E_{f}$ is the Young's modulus of the material, $\eta$
is the viscosity coefficient, $\nu$ is the Poisson's ratio, and $R$
is the radius of the spherical indenter [18]. The first term models
elastic deformation according to Hertzian contact mechanics, while
the second term accounts for viscoelastic damping.}

\textbf{2) Filter Design:} \pdfcomment[color=cyan,subject={},open=false,opacity=0.7]{User input: what was inferred from a given SINGLE force measurement?\\Claude comments: A single force sample drives one Bayesian update of an 8-dim EKF state: penetration depth $d$, velocity $\dot{d}$, local stiffness $k$, local damping $\lambda$, and their time derivatives. The contact-mechanics model (Eq.~6) relates $F$ to $(d, \dot{d}, E_f, \eta, \nu, R)$ --- the EKF inverts it online. $d$ itself is not directly measured; only $F$ is.}\sethlcolor{hlcyan}\hl{The main challenge lies in the fact that
$d$ is not directly measurable during robotic interaction. Therefore,
we use an EKF to estimate the system state, which includes both $d$
and $\dot{d}$, as well as the unknown tissue parameters. The system
is modelled using a discrete-time nonlinear state-space
representation, as detailed in [18]. The state vector is defined as}
\[
[\,d,\ k,\ \lambda,\ \dot{d},\ \dot{k},\ \dot{\lambda},\ \ddot{k},\
\ddot{\lambda}\,]^{\top},
\]
\sethlcolor{hlcyan}\hl{where $\kappa$ and $\lambda$ correspond to the local stiffness and
damping parameters, respectively. The filter assumes as input the
measured contact force and comprises the effective interaction mass.
Process and measurement noise are also included to account for model
uncertainties and sensor imperfections. This filtering approach
allows real-time estimation of the viscoelastic parameters even
during continuous contacts with indentation speeds up to 5 mm/s
[28], being suitable for dynamic exploration tasks.}

% ---------- Section III-E ----------
\subsection*{E.\ Target Distribution Update}

Once a sufficient number of new samples are acquired, the GPR model
is updated to refine the target stiffness distribution. Since GPR
inference has cubic time complexity with respect to the number of
samples, we avoid updating the model at every time step. Instead,
updates are performed at a fixed frequency (e.g., 1 Hz), allowing
batch accumulation of new data while keeping computation time
bounded.

% ---------- Section III-F ----------
\subsection*{F.\ Search Termination Criterion}

Different from existing approaches that rely on predefined stopping
conditions [15], or fixed time/step limits [16], we employ an
information-theoretic stopping criterion on $\alpha$, which is based
on the ergodic metric. The value of $\alpha$ governs the trade-off
between exploration and exploitation. Larger values of $\alpha$
correspond to an early exploration phase and favour earlier
termination, but may lead to incomplete spatial coverage. Conversely,
smaller values of $\alpha$ indicate that the exploration has largely
converged toward the target EID distribution, resulting in extended
exploitation phases, longer execution times, and only marginal
refinement of already identified regions. In practice, the stopping
threshold on $\alpha$ is selected within the exploitation-dominated
regime, such that further reductions correspond to diminishing
returns in reconstruction accuracy. This choice is supported by
empirical observations indicating saturation of the RMSE as $\alpha$
approaches a plateau, and reflects a practical balance between
coverage completeness, estimation accuracy, and exploration duration.
The advantage of using this metric lies in its ability to quantify
how thoroughly the area has been explored, without relying on prior
knowledge such as fixed time limits or the number and characteristics
of stiff regions that need to be found. This criterion is preferable
to a fixed-time limit, as it adapts the exploration duration to the
complexity of the stiffness distribution and the desired level of map
details, avoiding both premature termination and unnecessary probing.

% ---------- Section IV ----------
\section*{IV.\ Simulation Results}

\subsection*{A.\ Setup and Baselines}

The proposed search algorithm was evaluated on synthetic stiffness
distributions representative of biological soft tissues.

\newpage

% =========== Page 5: ergodic_page5.tex ===========
\setcounter{equation}{5}

% ---------- Fig. 2 placeholder ----------
\begin{center}
\fbox{\begin{minipage}[c][6cm][c]{0.95\textwidth}
\centering
\textbf{Fig. 2.}\ Synthetic elasticity distributions for simulated
scenarios, from 1 to 3 stiffer regions.\\[4pt]
\textcolor{gray}{\footnotesize
Three 45 mm $\times$ 45 mm square domains arranged side by side.
Left panel: a single stiff region; middle panel: two stiff regions;
right panel: three stiff regions. Axes $x_{1}, x_{2}$ in mm,
spanning 0--45. Colour scale: elastic modulus in kPa, ranging from
near-zero (soft background) up to approximately 180 kPa at the
stiff-inclusion peaks. Inclusions modelled as mixtures of Gaussian
components; background is near-uniform low stiffness.}
\end{minipage}}
\end{center}

\vspace{0.4em}

\textbf{Fig. 2.}\ Synthetic elasticity distributions for simulated
scenarios, from 1 to 3 stiffer regions. The colour scale refers to
the elastic modulus (in kPa).

\vspace{0.8em}

% ---------- Table I placeholder ----------
\begin{center}
\textbf{TABLE I}\\[2pt]
\textsc{Performance Metrics across 50 Runs with Noise-Free Elasticity
Measurements (w/o) and with White Noise $\mathcal{N}(0, 6.25\;
\mathrm{kPa}^{2})$ (w/). N\textsubscript{R} is the Number of Stiff
Regions, DR the Detection Rate, $l$ the Length of the Trajectory, and
RMSE is the Root Mean Square Error.}
\end{center}

\begin{center}
\footnotesize
\begin{tabular}{|c|c|c|c|c|}
\hline
$N_{R}$ & DR (\%) $\uparrow$ & $l$ [mm] & Time [s] $\downarrow$ &
RMSE [kPa] $\downarrow$ \\
\hline
\multicolumn{5}{|c|}{\textit{w/o Measurement Noise}} \\
\hline
1 & 100\% & 664 & 70 & 0.37 \\
2 & 100\% & 660 & 69 & 0.97 \\
3 & 99\%  & 665 & 70 & 1.96 \\
\hline
\multicolumn{5}{|c|}{\textit{w/ Measurement Noise}} \\
\hline
1 & 100\% & 633 & 66 & 1.57 \\
2 & 99\%  & 630 & 65 & 2.17 \\
3 & 97\%  & 580 & 60 & 3.14 \\
\hline
\end{tabular}
\end{center}

\vspace{0.8em}

% ---------- Prose continuation of Section IV-A ----------
They were generated as mixtures of Gaussian components with varying
locations, covariances, and amplitudes, modelling localised stiff
inclusions (over 100 kPa in pathological areas) in a softer
background (few kilopascal in healthy regions) [29], similarly to
related research [15]--[17]. Three scenarios in a square domain
(with side 45 mm) were considered: (i) a single stiff region,
(ii) two stiff regions, and (iii) three distinct stiff areas, whose
distributions are illustrated in Fig.~2. These regions were first
densely sampled with a grid-based approach to generate the ground
truth, used to compute the estimation and segmentation errors. The
simulation environment was implemented in Python, with GPR and BO
components built using the PyTorch framework, thus ensuring
efficient GPU-based inference and optimisation during online
execution. All simulations and experiments were executed on a mini
PC equipped with an Intel i7-13700HX CPU (24 cores), 64 GB RAM, and
an NVIDIA RTX 4070 GPU with 8 GB of memory. The proposed framework
runs in real time, with HEDAC trajectory updates and EKF-based
viscoelastic estimation both requiring less than 1 ms per control
loop iteration at 100 Hz. GPR updates are performed asynchronously
at a lower frequency (1 Hz), with a maximum observed update time of
105 ms. Elasticity samples were collected at a frequency of 4 Hz and
a maximum velocity of 1 cm/s. In this configuration, the maximum
distance between two consecutive sampled points equals half of the
indenter radius, which is set here to 2.5 mm. These simulations were
conducted to evaluate the performance of the proposed ergodic search
strategy (Sec.~IV-B) across the three distribution scenarios.

We compared ergodic search with two variants of Bayesian Optimisation
(BO), which is the standard baseline for robotic palpation
(Sec.~IV-C). BO models an expensive-to-evaluate objective function
$f(x)$ using a Gaussian Process (GP) and selects sampling locations
by maximising an acquisition function that balances exploration and
exploitation. Given data
$\mathcal{D} = (x_{i}, y_{i})_{i=1}^{n}$,
the GP provides a posterior predictive distribution
$f(x) \sim \mathcal{N}(\mu(x), \sigma^{2}(x))$. As an acquisition
function, the Expected Improvement (EI) [15], [16] is employed,
\[
a_{\mathrm{EI}}(x) =
(f_{\min} - \mu(x))\,\Phi(z) + \sigma(x)\,\phi(z),
\]
where $f_{\min}$ is the minimum observed value and
$z = \frac{f_{\min} - \mu(x)}{\sigma(x)}$. This variant is denoted
BO-EI. However, the standard point-to-point implementation of BO has
limitations for robotic palpation. Due to safety constraints,
especially in medical settings, the robot must follow smooth and
predictable trajectories and cannot move rapidly between distant
sampling points as typically generated by the BO. Moreover, each
probing action can take up to 5 s due to indentation and viscoelastic
estimation time. To ensure a fair comparison with the proposed
ergodic planner, we adapted BO to perform continuous palpation, which
characterises our ergodic search. This prevents BO from penalising
non-sampled regions during motion. Specifically, additional samples
are collected every 2.5 mm (half the indenter radius) along the path
between BO-selected targets, enabling BO to exploit intermediate
data, similarly to the ergodic approach (variant termed BO-EIS).

To assess the spatial accuracy of the reconstructed stiffness maps,
we applied a segmentation pipeline to identify distinct stiff
regions. Based on the elasticity estimates in Sec.~IV-B and IV-C,
KMeans clustering (SciPy) was used to separate stiffer regions from
the surrounding soft tissue. Unlike fixed thresholding approaches
[12], this method does not require manual parameter tuning and
adapts to the distribution of estimated stiffness values. The number
of detected regions was obtained by counting connected components in
the stiff cluster. In Sec.~IV-D, the segmented regions are converted
into polygonal boundaries using Shapely and compared against ground
truth using standard segmentation metrics.

% ---------- Section IV-B ----------
\subsection*{B.\ Robustness of Ergodic Exploration}

In this section, the performance of the ergodic controller is tested
with different numbers of stiffer regions. To assess the robustness
of the algorithm, for each map, the simulation is executed 50 times
with a random starting position, while three performance metrics
were considered: 1. the detection rate (DR), i.e., the percentage of
trials in which the algorithm correctly estimates the number and
position of the stiff regions present in the ground truth (higher is
better); 2. the elapsed time (lower is better); 3. the root mean
square error (RMSE, lower is better). For completeness, we report
also the trajectory length $l$. Each exploration terminates when the
ergodic metric $\alpha$ reaches 0.4 (set experimentally). Two
experimental conditions were examined: one involved noise-free
elasticity measurements and another incorporated noisy measurements
$\mathcal{N}(0, 6.25\;\mathrm{kPa}^{2})$, where the variance equals
the noise of the elasticity estimation in the worst case. The
results of the study are reported in Tab.~I. Overall, the algorithm
exhibited comparable performance in terms of DR and execution

\newpage

% =========== Page 6: ergodic_page6.tex ===========
\setcounter{equation}{5}

% ---------- Fig. 3 placeholder ----------
\begin{center}
\fbox{\begin{minipage}[c][8cm][c]{0.95\textwidth}
\centering
\textbf{Fig. 3.}\ End-effector tip trajectory during exploration of an
elasticity distribution with three stiff regions: planner comparison.\\[4pt]
\textcolor{gray}{\footnotesize
Five time snapshots in columns: $t = 14$ s, 28 s, 42 s, 56 s, 70 s.\\
Two rows: (top) Expected Information Density (EID) map; (bottom)
Estimated stiffness map.\\[3pt]
\textbf{(a)} Ergodic exploration with HEDAC (ours). The trajectory
smoothly covers the domain in proportion to EID; the cyan circle
marks trajectory start and the magenta square marks trajectory end.\\[2pt]
\textbf{(b)} BO-EI. Orange markers indicate BO-selected sampling
locations; red markers indicate points sampled along the path
between them.\\[2pt]
\textbf{(c)} BO-EIS. Same marker convention as (b); continuous
sampling augments BO's discrete selections.}
\end{minipage}}
\end{center}

\vspace{0.4em}

\textbf{Fig. 3.}\ End-effector tip trajectory during the exploration
of an elasticity distribution with three stiff regions: comparison of
different planners. The cyan circle indicates the beginning of the
trajectory, while the magenta square denotes the end of the
trajectory. In plots (b) and (c), the orange markers indicate the
locations selected by BO, whereas the red markers represent the
sampled points collected along the executed trajectory.

\vspace{0.8em}

% ---------- Section IV-B continuation ----------
time across the three scenarios, regardless of the presence of
measurement uncertainty. These results demonstrate the robustness of
the method with respect to both the number of stiff regions to be
detected and the influence of measurement noise. The RMSE, instead,
grows under both experimental conditions as the number of stiff
regions increases. This behaviour may be attributed to uncertainty
at the boundaries of the stiff regions, which are typically
under-sampled, and will be examined in greater detail in Sec.~IV-D.
As expected, we observe an increased RMSE in the second experimental
condition, attributable to the presence of the noise.

Figure~3(a) shows an example of the evolution of the agent trajectory
over time, when exploring the map with 3 stiff regions. All stiffer
regions are successfully detected within 40 s. Over time, the
ergodic metric $\alpha$ decreases from 0.79 at $t = 14$ s to 0.38 at
$t = 70$ s, reflecting improved coverage of the target distribution.
In parallel, the RMSE shows a substantial reduction, dropping from
13.75 kPa to 2.85 kPa. Intermediate values reinforce this trend: at
$t = 28$ s, $\alpha = 0.54$ and RMSE $= 11.19$ kPa; at $t = 42$ s,
$\alpha = 0.51$ and RMSE $= 3.7$ kPa; and at $t = 56$ s,
$\alpha = 0.43$ and RMSE $= 2.56$ kPa. This consistent evolution
confirms that as the exploration becomes more ergodic, the accuracy
of the reconstructed stiffness map improves accordingly.

% ---------- Section IV-C ----------
\subsection*{C.\ Comparison with BO Baselines}

In this section, we focus on the comparison of the results of our
ergodic approach with the two variants of BO described in Sec.~IV-A.
We consider the distribution with 3 stiff regions, which has proven
to be the most challenging for the ergodic exploration (see Fig.~3).
The simulation time is the sum of the simulated agent motion with a
velocity of 1 cm/s and the computation time required to determine
the next sampling location. Each simulation is conducted under two
conditions: one assuming noise-free elasticity measurements, and
another incorporating additive white noise
$\mathcal{N}(0, 6.25\;\mathrm{kPa}^{2})$, similar to Sec.~IV-B. Fixed
lenght of 400, 600, and 800 mm, were used to evaluate the performance
evolution over time, ensuring a fair comparison independently of the
method-specific stopping criteria. The average results of 100
simulations are listed in Tab.~II.

This comparison highlights that BO-EI consistently underperforms
compared to the other two methods, yielding the lowest DR and the
highest RMSE across all tested conditions. Both BO-EIS and the
ergodic strategy achieve comparable detection rates once sufficient
exploration has been performed, indicating successful identification
of the number and location of stiff regions. However, the ergodic
approach consistently provides lower reconstruction errors,
resulting in more accurate stiffness maps. At shorter travelled
distances (500 mm--650 mm), BO-EIS achieves higher detection rates,
reflecting faster initial identification of stiff regions. This
behaviour is expected given the exploitative nature of BO-based
strategies, which prioritise sampling high-uncertainty locations.
However, in clinical settings, once suspicious regions are
identified, accurate post-detection refinement becomes critical.
Therefore, it is more relevant to compare the approaches once
detection performance is maximised (800 mm, DR $\approx 99$--100\%),
and differences between the methods are primarily reflected in
reconstruction accuracy and execution duration. In this regime, the
ergodic strategy achieves lower RMSE while requiring substantially
less execution time than BO-EIS under the same travelled-distance
budget, both in noise-free and noisy conditions. These results
indicate that, in the post-detection phase, ergodic exploration
enables more efficient refinement of stiffness maps, yielding
improved accuracy per unit motion and reduced execution time under
equal motion cost. Compared to BO-based approaches, the ergodic
strategy also exhibits greater robustness to noise, maintaining a
consistent performance advantage across all evaluated conditions.
Overall, these findings confirm the benefit of ergodic exploration
in balancing detection reliability, reconstruction accuracy, and
execution efficiency once suspicious regions have been identified.

% ---------- Section IV-D ----------
\subsection*{D.\ Segmentation}

The spatial accuracy of the proposed method was evaluated based on
its ability to segment stiff regions within the reconstructed
stiffness maps and compared against the BO. To this end, we designed
two simulated stiffness distributions: the first one with two
inclusions, one irregularly shaped (Cluster 1) and one circular
(Cluster 2), while the second distribution consisted of a single
donut-shaped inclusion (Cluster 3), as

\newpage

% =========== Page 7: ergodic_page7.tex ===========
\setcounter{equation}{5}

% ---------- Table II ----------
\begin{center}
\textbf{TABLE II}\\[2pt]
\textsc{Comparison of performance metrics across 100 runs of Ergodic
Search and BO, with (w/) and without (w/o) noise. DR is the Detection
Rate, $t$ the Exploration Time, and RMSE the Root Mean Square Error.}
\end{center}

\begin{center}
\footnotesize
\begin{tabular}{|l|ccc|ccc|ccc|}
\hline
Traj.\ Length
  & \multicolumn{3}{c|}{500 mm}
  & \multicolumn{3}{c|}{650 mm}
  & \multicolumn{3}{c|}{800 mm} \\
\hline
Method & DR $\uparrow$ & $t\downarrow$ & RMSE$\downarrow$
       & DR $\uparrow$ & $t\downarrow$ & RMSE$\downarrow$
       & DR $\uparrow$ & $t\downarrow$ & RMSE$\downarrow$ \\
\hline
BO-EI (w/o)   & 39 & 59 & 12.09 & 73 & 88 & 7.98 & 92 & 124 & 5.47 \\
BO-EIS (w/o)  & 93 & 60 &  4.61 & 99 & 92 & 2.81 & 99 & 124 & 1.87 \\
Ergodic (w/o) & 78 & 41 &  4.69 & 92 & 62 & 2.74 & 99 &  84 & 1.42 \\
BO-EI (w/)    & 38 & 60 & 12.10 & 60 & 89 & 7.84 & 91 & 118 & 5.49 \\
BO-EIS (w/)   & 85 & 63 &  5.54 & 95 & 95 & 3.35 & 99 & 130 & 2.59 \\
Ergodic (w/)  & 79 & 41 &  5.87 & 98 & 62 & 3.03 & 99 &  84 & 2.16 \\
\hline
\end{tabular}
\end{center}

\vspace{0.6em}

% ---------- Table III ----------
\begin{center}
\textbf{TABLE III}\\[2pt]
\textsc{Average segmentation results across three clusters.}
\end{center}

\begin{center}
\footnotesize
\begin{tabular}{|c|ccc|ccc|ccc|}
\hline
 & \multicolumn{3}{c|}{Cluster 1}
 & \multicolumn{3}{c|}{Cluster 2}
 & \multicolumn{3}{c|}{Cluster 3} \\
\hline
 & BO-EI & BO-EIS & Erg. & BO-EI & BO-EIS & Erg. & BO-EI & BO-EIS & Erg. \\
\hline
Sens. & 0.561 & 0.950 & 0.975 & 0.657 & 0.997 & 0.971 & 0.746 & 0.904 & 0.978 \\
Spec. & 0.654 & 0.993 & 0.993 & 0.655 & 0.998 & 0.997 & 0.829 & 0.953 & 0.992 \\
\hline
\end{tabular}
\end{center}

\vspace{0.8em}

% ---------- Fig. 4 placeholder ----------
\begin{center}
\fbox{\begin{minipage}[c][4.5cm][c]{0.80\textwidth}
\centering
\textbf{Fig. 4.}\ Clustering and boundary extraction after exploration.\\[4pt]
\textcolor{gray}{\footnotesize
Original cluster shapes (in blue) overlaid against estimated shapes
(in red) for the three segmentation benchmarks. Cluster 1: irregular
inclusion; Cluster 2: circular inclusion; Cluster 3: donut-shaped
inclusion. The overlay visualises boundary fidelity relative to
ground truth.}
\end{minipage}}
\end{center}

\vspace{0.4em}

\textbf{Fig. 4.}\ Clustering and boundary extraction results after
the exploration phase. The original shapes (in blue) are shown
alongside the estimated ones (in red).

\vspace{0.8em}

% ---------- Section IV-D continuation ----------
shown in Fig.~4. A total of 50 simulations were conducted under
identical conditions, incorporating additive white noise
$\mathcal{N}(0, 6.25\;\mathrm{kPa}^{2})$. Each run was terminated
after 70 s to ensure consistent comparison across trials. For each
run, the estimated stiffness map was segmented using the clustering
and boundary extraction pipeline described earlier. The resulting
segmented areas were compared with ground truth using two standard
metrics [15], sensitivity and specificity, defined as:
\[
\text{1) Sensitivity} = \frac{\text{TP}}{\text{TP}+\text{FN}},
\quad
\text{2) Specificity} = \frac{\text{TN}}{\text{TN}+\text{FP}},
\]
where TP, FP, FN, and TN denote the number of pixels of true
positives, false positives, false negatives, and true negatives,
respectively. The metrics were computed separately for each of the
three ground-truth regions, and the final values were averaged
across the simulations. The results of the segmentation algorithm
(Tab.~III) demonstrate elevated and consistent specificity across
all three clusters for the two methods that sample continuously,
thereby confirming the model effectiveness in correctly excluding
background regions. Conversely, the ergodic search exhibited
superior sensitivity compared to both the BO with and without
sampling when clustering non-uniform shapes (Clusters 1--3). This
finding indicates that the ergodic search is more adept at
accurately delineating regions of interest. This is especially
relevant in medical contexts, where false negatives can lead to
missed detections of pathological regions. Consequently, achieving
high sensitivity is imperative to ensure clinically reliable
performance.

% ---------- Section V ----------
\section*{V.\ Experimental Results}

\pdfcomment[color=cyan,subject={},open=false,opacity=0.7]{User input: what exactly does their instrumentation measure in the ergodic paper?\\Claude comments: Directly measured: 6-axis contact force/torque (Bota SensOne at 1 kHz) + robot end-effector position (UR3e at 500 Hz). Stiffness, viscosity, penetration depth/velocity, and the stiffness map are all inferred downstream --- EKF over the DRM/Hertzian contact model (Eq.~6) gives per-point $E_f$ and $\eta$; GPR then interpolates to a global $\mu(x)$ surface.}\pdfcomment[color=cyan,subject={},open=false,opacity=0.7]{User input: how exactly does their instrumentation measure force?\\Claude comments: Paper names the commercial unit (Bota SensOne, 6-axis, 1 kHz) but does NOT describe the transduction mechanism --- treats the sensor as a black box. Deliberate architectural choice ('off-the-shelf ensures portability, cost-effectiveness'), explicitly contrasted with 'customised hardware integration'. Sharp architectural contrast with Chen, which uses a custom carbon-fibre piezoresistive sensor and cites its source paper~[35].}\pdfcomment[color=cyan,subject={},open=false,opacity=0.7]{User input: what is the basic description of their instrumentation which includes the force sensor in the ergodic paper?\\Claude comments: The canonical per-component rollcall: UR3e robot (500 Hz) $\to$ Bota SensOne 6-axis F/T sensor (1 kHz) $\to$ 3D-printed spherical indenter (5 mm radius) $\to$ ROS2 Humble middleware. The most complete instrumentation description in the paper; no rival passage comes close.}\pdfcomment[color=cyan,subject={},open=false,opacity=0.7]{User input: what part of the instrumentation does the PALPATION?\\Claude comments: The 3D-printed spherical indenter is the sole physical palpator --- the only component that contacts tissue. UR3e is the CARRIER (positions it), Bota SensOne is the OBSERVER (measures reaction force). Confirmed by Fig.~1 caption's explicit 'Indenter' label (p.~1) and Eq.~6's reliance on 5~mm radius $R$ as the Hertzian geometric parameter (p.~4). The one custom-fabricated element in an otherwise commercial-off-the-shelf instrumentation stack.}\sethlcolor{hlcyan}\hl{The validity of the simulation results in Sec.~IV was assessed in a
realistic setting by testing the proposed method on a manipulator
interacting with a physical sample (Fig.~1). An UR3e, position-%
controlled at 500 Hz, was used to execute the trajectories generated
by the Ergodic Control. The force at the end-effector was measured
by a 6-axis F/T sensor, a Bota SensOne, sampling at 1 kHz. A 3D-%
printed indenter with a spherical tip with 5 mm radius is attached
to the end-effector. The communication between the robot, the F/T
sensor, and the algorithms was implemented with ROS2 Humble.} To
simulate a biological tissue with a mass, we fabricated a silicone
sample consisting of a soft matrix made from Ecoflex-0030 and a
stiffer spherical inclusion (radius 1 cm) made from DragonSkin-30NV.
This configuration increases local stiffness by up to 90\%,
consistent with values reported in previous studies [15], [29]. The
ground truth is built with point-wise sampling on a grid with a
palpation every 2.5 mm in a square of 5 cm $\times$ 5 cm.

\pdfcomment[color=cyan,subject={},open=false,opacity=0.7]{User input: can their instrumentation palpate an object in the ergodic paper?\\Claude comments: Yes --- unambiguously. The UR3e + Bota SensOne + spherical indenter + ROS2 instrumentation successfully palpated a silicone phantom (Ecoflex matrix + DragonSkin-30NV spherical inclusion) in 5-out-of-5 trials from random initial positions. Hit rate 100\%, RMSE 9.76~kPa, sensitivity 0.92, specificity 0.98, with independent grid-based probing as ground truth.}\sethlcolor{hlcyan}\hl{We conducted 5 experiments on the physical setup, each starting from
a different random initial position. In all cases, the algorithm
successfully identified the location of the spherical intrusion,
demonstrating consistent target localisation. The reconstruction
accuracy (RMSE $= 9.76 \pm 2.97$ kPa) shows good agreement with the
reference stiffness distribution obtained from grid-based probing
used as ground truth. The method also achieved good classification
performance, with an average sensitivity of 0.92, and specificity of
0.98.} Compared to simulation results, the experimental metrics
confirm the high sensitivity and specificity, which are particularly
important for avoiding false negatives in medical applications.
Despite the experimental RMSE values being marginally higher than
those observed in simulation, the measurement uncertainty is still
below the 5\% (average absolute error $8.00 \pm 2.7$ kPa) and the
clustering performance appears largely unaffected. This difference
can be attributed to several factors, including the non-ideal
behaviour of real materials, noise in force measurements, slight
deviations from the assumed contact model, and uncertainties in
ground-truth stiffness calibration. Despite these challenges, the
observed trends are consistent with the simulation, confirming the
robustness of the method in real-world conditions. As shown in
Fig.~5, the trajectory of the end-effector adapts to the underlying
stiffness distribution: it concentrates more exploration in the
stiffer region on the left while still covering the rest of the map.

\newpage

% =========== Page 8: ergodic_page8.tex ===========
\setcounter{equation}{5}

% ---------- Fig. 5 placeholder ----------
\begin{center}
\fbox{\begin{minipage}[c][6cm][c]{0.90\textwidth}
\centering
\textbf{Fig. 5.}\ Estimated stiffness distribution from palpation of
the soft phantom.\\[4pt]
\textcolor{gray}{\footnotesize
50 mm $\times$ 50 mm domain in the $(x, y)$ plane. Colour scale:
estimated elastic modulus ranging from $\approx 102$ kPa (soft
background) up to $\approx 183$ kPa (stiff DragonSkin-30NV inclusion
shown in yellow). The end-effector trajectory during the search
motion is overlaid as a grayscale-gradient path: black marks the
beginning, white the end. The trajectory concentrates exploration in
the stiffer region on the left while still covering the rest of the
map, illustrating the adaptive focus of the EID-driven ergodic
planner.}
\end{minipage}}
\end{center}

\vspace{0.4em}

\textbf{Fig. 5.}\ Estimated stiffness distribution obtained through
palpation of the soft phantom. The stiffer region is shown in yellow.
The end-effector trajectory during the search motion is visualised
using a grayscale gradient, from black (beginning) to white (end). A
video of the experiment is provided in the Supplemental Materials
and is also available online at
\url{https://youtu.be/a3BUSxA7wXY}.

\vspace{0.8em}

% ---------- Section VI ----------
\section*{VI.\ Conclusion}

This work presents an autonomous robotic palpation framework for
real-time elastic stiffness mapping of soft tissues. By combining
force-based viscoelastic parameter estimation with ergodic trajectory
planning implemented via a heat-equation-driven controller, the
method enables continuous palpation without embedded tactile sensing.
Exploration is guided by an EID tailored to stiffness mapping,
allowing adaptive focus on diagnostically relevant regions.
Simulation results demonstrate improved stiffness reconstruction
accuracy, with reliable detection of multiple stiff regions and
lower RMSE and segmentation errors compared to BO, particularly
under noisy conditions. Experimental validation with a robotic arm
and a silicone phantom confirms the robustness of the approach,
showing consistent detection of stiff inclusions and good agreement
with ground-truth stiffness maps. Overall, the proposed ergodic
strategy effectively balances exploration and accurate estimation,
making it well suited for autonomous palpation as an assistive
diagnostic tool. Future work will extend the method to
three-dimensional mapping on non-flat surfaces and evaluate it on
anatomically realistic phantoms and ex vivo tissues.

% ---------- References ----------
\section*{References}

\begin{footnotesize}
\noindent
[1] R.~Tozzi, C.~K\"ohler, et al., ``Laparoscopic treatment of early
ovarian cancer: surgical and survival outcomes,'' Gynecologic
oncology, vol.~93, no.~1, pp.~199--203, 2004.\\[2pt]
[2] N.~Eisemann, S.~Bunk, et al., ``Nationwide real-world
implementation of AI for cancer detection in population-based
mammography screening,'' Nature Medicine 2025 31:3, vol.~31, no.~3,
pp.~917--924, 1 2025.\\[2pt]
[3] Q.~Dan, Z.~Xu, et al., ``Diagnostic performance of deep learning
in ultrasound diagnosis of breast cancer: a systematic review,'' npj
Precision Oncology 2024 8:1, vol.~8, no.~1, pp.~1--13, 1 2024.\\[2pt]
[4] A.~Attanasio, B.~Scaglioni, et al., ``Autonomy in Surgical
Robotics,'' Annual Review of Control, Robotics, and Autonomous
Systems, vol.~4, 2021.\\[2pt]
[5] L.~Zhang, Y.~Peng, et al., ``Innovative integration of automated
breast volume scan and ultrasound elastography for enhanced
differentiation of benign and malignant breast lesions,'' Scientific
Reports 2025 15:1, vol.~15, no.~1, pp.~1--8, 5 2025.\\[2pt]
[6] I.~Sack, ``Magnetic resonance elastography from fundamental
soft-tissue mechanics to diagnostic imaging,'' Nature Reviews
Physics, vol.~5, no.~1, 2023.\\[2pt]
[7] N.~Pacheco, C.~Henry, et al., ``Virtual palpation: A novel method
to identify the physical properties of tissue for laser surgery,''
in Hamlyn Symposium on Medical Robotics, 2025.\\[2pt]
[8] A.L.~Trejos, J.~Jayender, et al., ``Robot-assisted tactile
sensing for minimally invasive tumor localization,'' International
Journal of Robotics Research, vol.~28, no.~9, 2009.\\[2pt]
[9] N.~Enayati, E.~De~Momi, and G.~Ferrigno, ``Haptics in
robot-assisted surgery: Challenges and benefits,'' IEEE Reviews in
Biomedical Engineering, vol.~9, 2016.\\[2pt]
[10] S.~McKinley, A.~Garg, et al., ``An interchangeable surgical
instrument system with application to supervised automation of
multilateral tumor resection,'' in IEEE International Conference on
Automation Science and Engineering, 2016.\\[2pt]
[11] N.~Zevallos, R.~Arun~Srivatsan, et al., ``A Real-time Augmented
Reality Surgical System for Overlaying Stiffness Information,'' in
Robotics: Science and Systems, 2018.\\[2pt]
[12] A.~Garg, S.~Sen, et al., ``Tumor localization using automated
palpation with Gaussian Process Adaptive Sampling,'' in IEEE
International Conference on Automation Science and Engineering,
2016.\\[2pt]
[13] E.~Ayvali, R.A.~Srivatsan, et al., ``Using Bayesian optimization
to guide probing of a flexible environment for simultaneous
registration and stiffness mapping,'' in Proceedings - IEEE
International Conference on Robotics and Automation, 2016.\\[2pt]
[14] H.~Salman, E.~Ayvali, et al., ``Trajectory-Optimized Sensing for
Active Search of Tissue Abnormalities in Robotic Surgery,'' in
Proceedings - IEEE International Conference on Robotics and
Automation, 2018.\\[2pt]
[15] Y.~Yan and J.~Pan, ``Fast Localization and Segmentation of
Tissue Abnormalities by Autonomous Robotic Palpation,'' IEEE
Robotics and Automation Letters, vol.~6, no.~2, pp.~1707--1714,
4 2021.\\[2pt]
[16] E.~Ayvali, A.~Ansari, et al., ``Utility-Guided Palpation for
Locating Tissue Abnormalities,'' IEEE Robotics and Automation
Letters, vol.~2, no.~2, 2017.\\[2pt]
[17] K.A.~Nichols and A.M.~Okamura, ``Methods to Segment Hard
Inclusions in Soft Tissue During Autonomous Robotic Palpation,''
IEEE Transactions on Robotics, vol.~31, no.~2, 2015.\\[2pt]
[18] L.~Beber, E.~Lamon, et al., ``Towards Robotised Palpation for
Cancer Detection through Online Tissue Viscoelastic Characterisation
with a Collaborative Robotic Arm,'' in IEEE International Conference
on Intelligent Robots and Systems, 2024, pp.~2380--2386.\\[2pt]
[19] S.~Ivic, B.~Crnkovic, and I.~Mezic, ``Ergodicity-Based
Cooperative Multiagent Area Coverage via a Potential Field,'' IEEE
Transactions on Cybernetics, vol.~47, no.~8, 2017.\\[2pt]
[20] C.E.~Rasmussen and C.K.~Williams, Gaussian Processes for
Machine Learning. The MIT Press, 2006.\\[2pt]
[21] P.~Chalasani, L.~Wang, et al., ``Concurrent nonparametric
estimation of organ geometry and tissue stiffness using continuous
adaptive palpation,'' in Proceedings - IEEE International Conference
on Robotics and Automation, vol.~2016-June, 2016.\\[2pt]
[22] ---, ``Preliminary evaluation of an online estimation method
for organ geometry and tissue stiffness,'' IEEE Robotics and
Automation Letters, vol.~3, no.~3, 2018.\\[2pt]
[23] G.~Mathew and I.~Mezi, ``Metrics for ergodicity and design of
ergodic dynamics for multi-agent systems,'' Physica D: Nonlinear
Phenomena, vol.~240, no.~4-5, pp.~432--442, 2 2011.\\[2pt]
[24] L.M.~Miller, Y.~Silverman, et al., ``Ergodic Exploration of
Distributed Information,'' IEEE Transactions on Robotics, vol.~32,
no.~1, pp.~36--52, 2 2016.\\[2pt]
[25] L.M.~Miller and T.D.~Murphey, ``Trajectory optimization for
continuous ergodic exploration,'' in Proceedings of the American
Control Conference, 2013.\\[2pt]
[26] C.~Bilaloglu, T.~L\"ow, and S.~Calinon, ``Whole-body ergodic
exploration with a manipulator using diffusion,'' IEEE Robotics and
Automation Letters, vol.~8, no.~12, pp.~8581--8587, 2023.\\[2pt]
[27] V.L.~Popov and M.~He\ss, Method of dimensionality reduction in
contact mechanics and friction. Springer, 2015.\\[2pt]
[28] L.~Beber, E.~Lamon, et al., ``Force-based viscosity and
elasticity measurements for material biomechanical characterisation
with a collaborative robotic arm,'' IEEE Transactions on
Instrumentation and Measurement, pp.~1--1, 2025.\\[2pt]
[29] P.N.~Wells and H.-D.~Liang, ``Medical ultrasound: imaging of
soft tissue strain and elasticity,'' Journal of the Royal Society
Interface, vol.~8, no.~64, pp.~1521--1549, 2011.
\end{footnotesize}

\end{document}
